% ============================================================
%  ТЕСТИРОВАНИЕ И ОЦЕНКА
% ============================================================

\section{Тестирование и оценка}

\subsection{Общий подход к тестированию}

Инструмент реализован как конвейер из шести взаимозависимых модулей,
каждый из которых находится на критическом пути итогового отчёта.
Их взаимосвязь показана на рисунке~\ref{fig:pipeline}.

\begin{figure}[h]
  \centering
  \includegraphics[width=\textwidth]{pipeline.pdf}
  \caption{Конвейер обработки данных}
  \label{fig:pipeline}
\end{figure}

Поскольку ошибка в любом из звеньев цепи может проявиться только на
выходе, применяются два уровня тестирования.
\textbf{Модульные тесты} (41~тест, 6~наборов) проверяют каждый
компонент изолированно~--- это позволяет немедленно локализовать
источник ошибки.
\textbf{Интеграционные тесты} (6~тестов) запускают полный конвейер
на реальном тестовом файле и сверяют итоговый план миграции с
эталонной разметкой.

Все тесты реализованы на~C++ с использованием собственного
фреймворка \texttt{TestSuite\,/\,TestRunner}, не требующего
сторонних зависимостей.
Итого выполняется \textbf{47 тестов}, все проходят успешно.
Классификация приведена в таблице~\ref{tab:testing_levels}.

\begin{longtable}{|p{4.8cm}|p{5.0cm}|p{5.2cm}|}
  \caption{Классификация тестов системы}
  \label{tab:testing_levels} \\
  \hline
  Набор тестов (кол-во) & Что проверяется & Ключевые тест"=кейсы \\ \hline
  \endfirsthead
  \multicolumn{3}{r}{\textit{Продолжение таблицы \thetable}} \\ \hline
  Набор тестов (кол-во) & Что проверяется & Ключевые тест"=кейсы \\ \hline
  \endhead \hline \endlastfoot

  \texttt{VulnDatabase} (9)
    & Загрузка JSON, поиск по имени и псевдониму, слияние с
      дополнительной базой, наличие ссылок на FIPS и ТК\,26
    & \texttt{find\_alias}, \texttt{replacements\_present},
      \texttt{merge\_extra\_db} \\ \hline

  \texttt{ProjectScanner} (5)
    & Обход файловой системы, классификация по категориям,
      статистика языков, метаданные CBOM, плагин
    & \texttt{cpp\_files\_detected},
      \texttt{cbom\_metadata\_json},
      \texttt{custom\_plugin} \\ \hline

  \texttt{RegexAnalyzer} (8)
    & Обнаружение функций, атрибуты находки (строка, библиотека,
      ссылка NIST), отсутствие срабатываний на AES
    & \texttt{finds\_rsa\_in\_fixture},
      \texttt{no\_false\_positives\_for\_safe\_code},
      \texttt{finding\_has\_nist\_reference} \\ \hline

  \texttt{ASTAnalyzer} (7)
    & Обнаружение уязвимостей, извлечение контекста функции и
      класса, исключение системных заголовков
    & \texttt{context\_function\_extracted},
      \texttt{context\_class\_extracted},
      \texttt{no\_system\_header\_findings} \\ \hline

  \texttt{RiskAssessor} (7)
    & Влияние контекста на итоговую оценку, сортировка плана
      миграции, наличие рекомендаций и ссылок на стандарты
    & \texttt{network\_file\_increases\_risk},
      \texttt{test\_file\_reduces\_risk},
      \texttt{replacement\_info\_present} \\ \hline

  \texttt{MetricsCalculator} (5)
    & Вычисление TP, FP, FN с допуском строки;
      precision, recall, F1; сериализация в JSON
    & \texttt{perfect\_match}, \texttt{line\_tolerance},
      \texttt{json\_serialization} \\ \hline

  \texttt{FullPipeline} (6)
    & Полный конвейер в режимах regex и AST;
      precision/recall\,$\geq$\,50\,\%;
      наличие ссылок ТК\,26 и FIPS; сериализация отчёта
    & \texttt{regex\_pipeline\_end\_to\_end},
      \texttt{regex\_precision\_recall\_gt50pct},
      \texttt{migration\_plan\_has\_fips\_references} \\

\end{longtable}


% -----------------------------------------------------------
\subsection{Тестовые данные}

Основным тестовым файлом является
\texttt{tests/\allowbreak fixtures/\allowbreak sample\_vulnerable.cpp}.
Он организован в виде восьми тематических секций~(A--H), каждая из
которых моделирует отдельный реальный сценарий применения
криптографии, и завершающего блока \texttt{SAFE}~--- фрагментов,
которые \emph{не должны} давать срабатываний.
Такое разделение позволяет измерять precision и recall раздельно
по категориям уязвимостей.

Структура тестового файла представлена в
таблице~\ref{tab:fixture_sections}.

\begin{longtable}{|p{2.0cm}|p{4.8cm}|p{8.2cm}|}
  \caption{Секции тестового файла \texttt{sample\_vulnerable.cpp}}
  \label{tab:fixture_sections} \\
  \hline
  Секция & Класс"=контейнер & Описание и ожидаемые находки \\ \hline
  \endfirsthead
  \multicolumn{3}{r}{\textit{Продолжение таблицы \thetable}} \\ \hline
  Секция & Класс"=контейнер & Описание и ожидаемые находки \\ \hline
  \endhead \hline \endlastfoot

  A.~\texttt{tls}
    & \texttt{crypto::\allowbreak tls::\allowbreak TLSHandshake}
    & RSA"=ключ сервера: \texttt{RSA\_new}, \texttt{RSA\_generate\_key\_ex};
      EC"=ключ и ECDH: \texttt{EC\_KEY\_generate\_key},
      \texttt{ECDH\_compute\_key}; подпись~ECDSA \\ \hline

  B.~\texttt{storage}
    & \texttt{crypto::\allowbreak storage::\allowbreak SecureStorage}
    & RSA"=шифрование: \texttt{RSA\_public\_encrypt};
      подпись~RSA; хеширование~SHA"=1 \\ \hline

  C.~DH
    & функция \texttt{setup\_dh}
    & Параметры Диффи"=Хеллмана:
      \texttt{DH\_new}, \texttt{DH\_generate\_parameters\_ex} \\ \hline

  D.~MD5
    & функция \texttt{md5sum}
    & Устаревшая хеш"=функция: \texttt{MD5} \\ \hline

  E.~\texttt{evp}
    & \texttt{crypto::\allowbreak evp::\allowbreak EvpOperations}
    & EVP"=API высокого уровня: \texttt{EVP\_PKEY\_CTX\_new\_from\_name}
      с классическим аргументом \texttt{"RSA"}; \texttt{EVP\_PKEY\_CTX\_new};
      полная цепочка подписи и верификации
      \texttt{EVP\_DigestSign*}/\texttt{EVP\_DigestVerify*};
      ECDH"=согласование через \texttt{EVP\_PKEY\_derive*} \\ \hline

  F.~\texttt{dsa}
    & \texttt{crypto::\allowbreak dsa\_legacy::\allowbreak DsaOps}
    & Явные DSA"=операции: \texttt{DSA\_new}, \texttt{DSA\_sign};
      проверяет, что анализатор не смешивает~DSA и~ECDSA \\ \hline

  G.~\texttt{ec\_ops}
    & \texttt{crypto::\allowbreak ec\_ops::\allowbreak ECOperations}
    & Низкоуровневые EC"=операции:
      \texttt{EC\_GROUP\_new\_by\_curve\_name},
      \texttt{EC\_KEY\_new}, \texttt{EC\_KEY\_generate\_key} \\ \hline

  H.~\texttt{bn\_ops}
    & \texttt{crypto::\allowbreak bn\_ops::\allowbreak BNOperations}
    & BN"=операции как индикатор ручной реализации RSA/DH:
      \texttt{BN\_generate\_prime\_ex}, \texttt{BN\_mod\_exp} \\ \hline

  SAFE
    & функции \texttt{safe\_*}
    & AES"=256"=GCM через \texttt{EVP\_EncryptInit\_ex}:
      симметричный алгоритм, отсутствует в базе уязвимых функций. \\
    & & \texttt{EVP\_PKEY\_CTX\_new\_from\_name} с аргументом
      \texttt{"ML"=KEM"=768"}: постквантовый алгоритм,
      подавляется AST"=анализатором через отслеживание провенанса
      аргумента \\

\end{longtable}

Полная эталонная разметка из \textbf{32 записей} хранится в файле
\texttt{tests/\allowbreak fixtures/\allowbreak ground\_truth.json}.
Каждая запись включает: имя ожидаемой функции, номер строки и
\texttt{line\_tolerance}~(3--5~строк)~--- допуск на расхождение
номеров строк между AST и regex"=режимами.
Перечень всех ожидаемых находок приведён в
приложении~\ref{app:ground_truth}.

Листинг~\ref{lst:gt-entry-example} показывает типичную структуру
записи в \texttt{ground\_truth.json}.

\begin{itmolisting}[JSON]{Структура записи эталонной разметки}{lst:gt-entry-example}
{
  "function_name": "EVP_PKEY_CTX_new_from_name",
  "file_path": "tests/fixtures/sample_vulnerable.cpp",
  "line_number": 121,
  "line_tolerance": 3,
  "notes": "Conditional: 'RSA' arg makes it vulnerable; 'ML-KEM-768' must NOT trigger"
}
\end{itmolisting}

Поле \texttt{notes} носит документальный характер и не используется
в расчёте метрик.
Допуск \texttt{line\_tolerance} необходим потому, что AST"=анализатор
привязывает находку к точной строке вызова, тогда как regex"=анализатор
может зафиксировать строку начала многострочного выражения.


% -----------------------------------------------------------
\subsection{Модульные тесты}

\subsubsection{База данных уязвимых функций (\texttt{VulnDatabase})}

\textbf{Назначение и структура.}
Класс \texttt{VulnDatabase} загружает JSON"=файл
\texttt{data/\allowbreak vulnerable\_functions.json},
строит инвертированный индекс по именам и псевдонимам
и предоставляет метод \texttt{find\_by\_name(name)} с нечувствительным
к регистру поиском.
Каждая запись в базе представлена структурой \texttt{VulnerableFunction}:

\begin{itmolisting}[C++]{Структура записи базы данных уязвимых функций}{lst:vf-struct}
struct VulnerableFunction {
    std::string id, name;          // каноническое имя (напр., RSA_generate_key_ex)
    std::string category;          // asymmetric_key_generation | key_exchange | ...
    std::string algorithm;         // RSA | ECDSA | DH | SHA-1 | ...
    std::string quantum_vulnerability; // high | medium | conditional
    std::string nist_reference, tc26_reference;
    double risk_score = 0.0;       // базовая оценка риска [0..10]
    std::vector<std::string> aliases;  // альтернативные имена (RSA_generate_key)
    std::vector<std::string> patterns; // regex-паттерны для быстрого поиска
    std::vector<ReplacementInfo> replacements; // рекомендованные замены
};
\end{itmolisting}

Механизм псевдонимов позволяет одной записи базы описывать несколько
вариантов API: например, устаревший \texttt{RSA\_generate\_key} и
актуальный \texttt{RSA\_generate\_key\_ex} объединены в одну запись
с каноническим именем \texttt{RSA\_generate\_key\_ex}.
Метод \texttt{merge\_json} обеспечивает расширяемость: организации
могут добавлять собственные записи для корпоративных криптографических
обёрток без изменения основной базы.

\textbf{Тесты.}
Тест \texttt{find\_alias} (листинг~\ref{lst:db-find-alias}) проверяет
ключевое свойство: поиск по псевдониму должен возвращать запись
с каноническим именем.
Это важно, поскольку анализаторы обнаруживают фактически вызванную
функцию, но риск и рекомендации хранятся при каноническом имени.

\begin{itmolisting}[C++]{Тест поиска по псевдониму функции}{lst:db-find-alias}
TEST(VulnDatabase, find_alias) {
    pqc::VulnDatabase db;
    db.load("data/vulnerable_functions.json");

    // RSA_generate_key — устаревшее имя, псевдоним канонического
    // RSA_generate_key_ex. Оба имени должны возвращать одну запись.
    auto opt = db.find_by_name("RSA_generate_key");

    ASSERT_TRUE(opt.has_value());
    ASSERT_EQ(opt->name, "RSA_generate_key_ex"); // возвращается каноническое имя
    ASSERT_EQ(opt->algorithm, "RSA");
}
\end{itmolisting}

Тест \texttt{replacements\_present} проверяет, что поле
\texttt{standard} в первой рекомендации по миграции содержит ссылку
на FIPS~--- это обязательный атрибут для формирования раздела
\texttt{replacement} в плане миграции:

\begin{itmolisting}[C++]{Тест наличия рекомендации по миграции}{lst:db-replacement}
TEST(VulnDatabase, replacements_present) {
    pqc::VulnDatabase db;
    db.load("data/vulnerable_functions.json");

    auto opt = db.find_by_name("RSA_generate_key_ex");

    ASSERT_TRUE(opt.has_value());
    ASSERT_NOT_EMPTY(opt->replacements);            // запись содержит рекомендацию
    ASSERT_CONTAINS(opt->replacements[0].standard, "FIPS"); // ссылка на стандарт
}
\end{itmolisting}

Полный листинг теста \texttt{merge\_extra\_db}, демонстрирующего
добавление корпоративной функции"=обёртки, приведён в
приложении~\ref{app:tests_db}.

% -----------------------------------------------------------
\subsubsection{Инвентаризация проекта (\texttt{ProjectScanner})}

\textbf{Назначение и структура.}
Класс \texttt{ProjectScanner} рекурсивно обходит директорию проекта
и строит объект \texttt{ProjectInventory}~--- центральный
носитель метаданных, который передаётся всем анализаторам.

\begin{itmolisting}[C++]{Ключевые структуры инвентаризации проекта}{lst:inventory-struct}
struct FileEntry {
    std::filesystem::path path;
    FileCategory category; // SOURCE | HEADER | BUILD | TEST | CERTIFICATE | ...
    std::string language;  // "cpp", "cmake", "x509"
    int line_count = 0;
};

struct ProjectInventory {
    std::string project_name;
    std::vector<FileEntry> files;
    std::map<std::string, int> language_stats; // {"cpp": 12, "cmake": 2, ...}

    // Вспомогательный метод: возвращает только исходные файлы (category == SOURCE)
    std::vector<const FileEntry*> get_source_files() const;

    // Формирует метаданные CBOM (имя проекта, статистика файлов, языки)
    nlohmann::json to_cbom_metadata() const;
};
\end{itmolisting}

Классификация файлов реализована через плагины
\texttt{ILanguagePlugin}: встроенные плагины обрабатывают C/C++,
CMake и X.509"=сертификаты; пользовательский плагин можно
зарегистрировать через \texttt{register\_plugin()}.

\textbf{Тесты.}
Тест \texttt{cpp\_files\_detected} проверяет, что после сканирования
директории \texttt{tests/\allowbreak fixtures} в инвентаре присутствуют
файлы \texttt{.cpp} в категории \texttt{SOURCE}~--- это предусловие
для корректной работы обоих анализаторов:

\begin{itmolisting}[C++]{Тест обнаружения C++ файлов при инвентаризации}{lst:scanner-cpp}
TEST(ProjectScanner, cpp_files_detected) {
    pqc::ProjectScanner scanner;
    auto inv = scanner.scan("tests/fixtures");

    // get_source_files() — фильтр по FileCategory::SOURCE
    auto src = inv.get_source_files();

    ASSERT_GT((int)src.size(), 0);
    bool found_cpp = false;
    for (auto* f : src)
        if (f->path.extension() == ".cpp") { found_cpp = true; break; }
    ASSERT_TRUE(found_cpp); // тестовый файл должен быть классифицирован как SOURCE
}
\end{itmolisting}

% -----------------------------------------------------------
\subsubsection{Regex"=анализатор (\texttt{RegexAnalyzer})}

\textbf{Назначение и структура.}
\texttt{RegexAnalyzer} при инициализации компилирует паттерны из базы
данных в объекты \texttt{std::regex} и хранит их с указателями на
соответствующие записи \texttt{VulnerableFunction}.
При анализе файла анализатор построчно применяет все паттерны
к строке, предварительно очищенной от блочных комментариев.
Результатом является набор объектов \texttt{Finding}.

Структура \texttt{Finding} служит универсальным носителем
информации об уязвимости и передаётся между всеми последующими
модулями конвейера:

\begin{itmolisting}[C++]{Структура результата обнаружения уязвимости}{lst:finding-struct}
struct Finding {
    std::string file_path;
    std::string function_name;   // фактически вызванная функция (может быть псевдонимом)
    std::string library;         // "openssl", "cryptopp", "gost_openssl"
    std::string algorithm;       // "RSA", "ECDH", "SHA-1", ...
    std::string category;        // "asymmetric_key_generation" | "digital_signature" | ...
    std::string quantum_vulnerability; // "high" | "medium" | "conditional"
    std::string context_function;  // объемлющая функция (заполняется AST-режимом)
    std::string context_class;     // объемлющий класс/структура
    std::string nist_reference, tc26_reference;
    int line_number = 0;
    double base_risk_score = 0.0;
    std::vector<std::string> arguments; // аргументы вызова (для conditional-проверок)
};
\end{itmolisting}

\textbf{Тесты.}
Тест \texttt{no\_false\_positives\_for\_safe\_code}
(листинг~\ref{lst:regex-safe}) является критическим тестом
корректности: он проверяет, что симметричный вызов
\texttt{EVP\_EncryptInit\_ex} в блоке \texttt{SAFE} не порождает
находок, поскольку функция отсутствует в базе уязвимых функций
(AES"=256"=GCM~--- симметричный режим, устойчивый при длине ключа
256~бит даже с учётом алгоритма Гровера).

\begin{itmolisting}[C++]{Тест отсутствия ложных срабатываний на симметричный AES-код}{lst:regex-safe}
TEST(RegexAnalyzer, no_false_positives_for_safe_code) {
    pqc::VulnDatabase db;
    db.load("data/vulnerable_functions.json");
    pqc::RegexAnalyzer ra(db);
    auto findings = ra.analyze_file(
        "tests/fixtures/sample_vulnerable.cpp"
    );

    // EVP_EncryptInit_ex с AES-256-GCM не должен входить в базу уязвимых функций.
    // Перебор всех находок гарантирует отсутствие данной функции в любом месте.
    for (auto& f : findings)
        ASSERT_TRUE(f.function_name != "EVP_EncryptInit_ex");
}
\end{itmolisting}

Тест \texttt{finding\_has\_nist\_reference} перебирает все находки
и проверяет, что каждая несёт непустую строку ссылки на стандарт
NIST~--- поле обязательно для раздела \texttt{nist\_compliance} в
JSON"=отчёте.

% -----------------------------------------------------------
\subsubsection{AST"=анализатор (\texttt{ASTAnalyzer})}

\textbf{Назначение и структура.}
\texttt{ASTAnalyzer} использует libclang для построения
дерева разбора исходного файла.
Для каждого узла"=вызова (\texttt{CXCursor\_CallExpr}) анализатор:
\begin{enumerate}
  \item определяет имя вызываемой функции через
        \texttt{clang\_getCursorReferenced};
  \item выполняет поиск в базе \texttt{VulnDatabase} по этому имени;
  \item при обнаружении условно уязвимой функции
        (\texttt{quantum\_vulnerability = "conditional"}) прослеживает
        провенанс аргументов через \texttt{resolve\_expr}, обходя
        цепочки инициализации переменных назад до литерала или вызова;
  \item подавляет находку, если аргумент содержит имя постквантового
        алгоритма (ml"=kem, ml"=dsa, slh"=dsa и~др.);
  \item извлекает контекст: объемлющую функцию, класс и пространство
        имён через обход родительских курсоров.
\end{enumerate}
При отсутствии libclang анализатор автоматически переключается на
токенизирующий режим с пониженной точностью извлечения контекста.

\textbf{Тесты.}
Тест \texttt{context\_class\_extracted} (листинг~\ref{lst:ast-ctx-class})
проверяет, что AST"=анализатор действительно извлекает объемлющий
класс.
Это важно, поскольку поле \texttt{context\_class} используется
\texttt{RiskAssessor} при расчёте контекстных множителей и
отображается в отчёте для понимания местоположения уязвимости.

\begin{itmolisting}[C++]{Тест извлечения объемлющего класса AST-анализатором}{lst:ast-ctx-class}
TEST(ASTAnalyzer, context_class_extracted) {
    pqc::VulnDatabase db;
    db.load("data/vulnerable_functions.json");
    pqc::ASTAnalyzer aa(db);
    auto findings = aa.analyze_file(
        "tests/fixtures/sample_vulnerable.cpp"
    );

    // Тестовый файл содержит уязвимые вызовы внутри нескольких классов
    // (TLSHandshake, SecureStorage, EvpOperations и др.).
    // Хотя бы для одной находки context_class должно быть непустым.
    bool has_class = false;
    for (auto& f : findings)
        if (!f.context_class.empty()) { has_class = true; break; }
    ASSERT_TRUE(has_class);
}
\end{itmolisting}

Тест \texttt{no\_system\_header\_findings} проверяет критическое
ограничение: все находки должны указывать на файл пользователя, а
не на системные заголовки OpenSSL.
Без этой проверки анализатор мог бы ошибочно репортировать
внутренние объявления библиотечных функций как уязвимые вызовы
в проекте.

% -----------------------------------------------------------
\subsubsection{Оценщик риска (\texttt{RiskAssessor})}

\textbf{Назначение и структура.}
\texttt{RiskAssessor} принимает список объектов \texttt{Finding},
вычисляет для каждого итоговую оценку риска и формирует
\texttt{RiskReport} с отсортированным планом миграции.

Итоговая оценка вычисляется как
$\text{final\_score} = \min(10{,}0,\;\text{base\_score} \times m)$,
где мультипликатор~$m$ формируется из шести независимых факторов,
описанных в структуре \texttt{ContextFactors}:

\begin{itmolisting}[C++]{Контекстные факторы и их влияние на оценку риска}{lst:context-factors}
struct ContextFactors {
    bool is_network_facing  = false; // файл связан с сетью  → ×1.3
    bool is_persistent_data = false; // файл связан с хранилищем → ×1.25
    bool is_key_material    = false; // генерация/обмен/подпись → ×1.2
    bool is_in_loop         = false; // вызов в цикле → ×1.1
    bool is_test_code       = false; // путь содержит /test → ×0.7
    int  call_frequency     = 1;     // частота > 5 → ×1.1
    std::string data_classification; // "key_material" | "session" | ...
};
\end{itmolisting}

На выходе \texttt{RiskAssessor::assess} возвращает \texttt{RiskReport}~---
план миграции, отсортированный по убыванию \texttt{final\_score},
со счётчиками по приоритетам и сводкой по файлам:

\begin{itmolisting}[C++]{Структуры выходных данных оценщика риска}{lst:risk-report-struct}
struct MigrationAction {
    std::string file_path, current_function, current_library;
    std::string replacement_function, replacement_standard; // напр., "FIPS 203"
    std::string nist_reference, tc26_note;
    RiskScore   risk;          // base_score, context_multiplier, final_score
    MigrationPriority priority; // CRITICAL | HIGH | MEDIUM | LOW
    int migration_order;        // порядковый номер в плане (1 — первоочередной)
};

struct RiskReport {
    std::vector<MigrationAction> migration_plan; // отсортировано по убыванию риска
    std::map<std::string, FileRiskSummary> file_summaries; // сводка по файлам
    double overall_risk_score; // средневзвешенная оценка по топ-10 + сертификаты
    std::string nist_readiness; // "not-started" | "in-progress" | "near-compliant" | "compliant"
};
\end{itmolisting}

\textbf{Тесты.}
Тест \texttt{migration\_order\_sorted\_by\_risk}
(листинг~\ref{lst:risk-order}) намеренно передаёт две находки
в обратном порядке риска и проверяет, что план миграции
упорядочен правильно.

\begin{itmolisting}[C++]{Тест сортировки плана миграции по убыванию риска}{lst:risk-order}
TEST(RiskAssessor, migration_order_sorted_by_risk) {
    pqc::VulnDatabase db;
    db.load("data/vulnerable_functions.json");
    pqc::RiskAssessor ra(db);
    pqc::ProjectInventory inv;

    // SHA1 (hash_function, base_risk=6.0) намеренно передаётся первым,
    // RSA_generate_key_ex (asymmetric_key_gen, base_risk=10.0) — вторым.
    pqc::Finding f1, f2;
    f1.function_name = "SHA1";
    f1.base_risk_score = 6.0; f1.file_path = "src/util.cpp";
    f1.category = "hash_function";

    f2.function_name = "RSA_generate_key_ex";
    f2.base_risk_score = 10.0; f2.file_path = "src/net.cpp";
    f2.category = "asymmetric_key_generation";

    auto rr = ra.assess({f1, f2}, inv); // неупорядоченный ввод

    // После оценки план миграции должен быть отсортирован по убыванию риска:
    // migration_plan[0] = RSA (выше риск), migration_plan[1] = SHA1.
    ASSERT_EQ((int)rr.migration_plan.size(), 2);
    ASSERT_GE(rr.migration_plan[0].risk.final_score,
              rr.migration_plan[1].risk.final_score);
}
\end{itmolisting}

Тест \texttt{test\_file\_reduces\_risk} создаёт находку с путём
\texttt{tests/\allowbreak test\_tls.cpp} и проверяет, что итоговая
оценка ниже базовой: \texttt{RiskAssessor} распознаёт путь как
тестовый файл и применяет коэффициент~$\times0{,}7$.

% -----------------------------------------------------------
\subsubsection{Модуль метрик (\texttt{MetricsCalculator})}

\textbf{Назначение и структура.}
\texttt{MetricsCalculator} сопоставляет список объектов
\texttt{Finding} с эталонной разметкой (\texttt{GroundTruth}),
используя двойной критерий совпадения: имя функции и номер строки
в пределах допуска \texttt{line\_tolerance}.
Результат~--- объект \texttt{AnalysisMetrics} с полями
\texttt{true\_positives}, \texttt{false\_positives},
\texttt{false\_negatives} и методами \texttt{precision()},
\texttt{recall()}, \texttt{f1\_score()}.

Тест \texttt{line\_tolerance} (листинг~\ref{lst:metrics-tol})
проверяет корректность допуска: находка на строке~45 должна
засчитываться как TP при эталонном значении~42 и допуске~5.

\begin{itmolisting}[C++]{Тест допуска при сопоставлении строк}{lst:metrics-tol}
TEST(MetricsCalculator, line_tolerance) {
    pqc::MetricsCalculator calc;

    pqc::Finding f;
    f.function_name = "RSA_sign";
    f.line_number   = 45; // AST привязывает к строке вызова внутри выражения

    pqc::GroundTruth gt;
    gt.function_name  = "RSA_sign";
    gt.line_number    = 42; // regex фиксирует начало выражения
    gt.line_tolerance = 5;  // допуск перекрывает оба варианта

    auto m = calc.compute({f}, {gt});

    // |45 - 42| = 3 <= 5, поэтому совпадение засчитывается как TP.
    ASSERT_EQ(m.true_positives,  1);
    ASSERT_EQ(m.false_positives, 0);
    ASSERT_EQ(m.false_negatives, 0);
}
\end{itmolisting}

% -----------------------------------------------------------
\subsection{Интеграционные тесты}

Интеграционные тесты проверяют полный конвейер обработки
на директории \texttt{tests/\allowbreak fixtures}.
Последовательность вызовов в каждом тесте отражает реальную
цепочку исполнения инструмента:

\[
  \texttt{VulnDatabase}
  \;\to\; \texttt{ProjectScanner}
  \;\to\; \texttt{Analyzer}
  \;\to\; \texttt{RiskAssessor}
  \;\to\; \texttt{RiskReport}
\]

Тест \texttt{regex\_pipeline\_end\_to\_end}
(листинг~\ref{lst:itest-regex}) последовательно выполняет все
этапы и проверяет непустоту плана миграции, ненулевую общую
оценку риска и наличие строки \texttt{nist\_readiness}.

\begin{itmolisting}[C++]{Интеграционный тест полного конвейера в режиме regex}{lst:itest-regex}
TEST(FullPipeline, regex_pipeline_end_to_end) {
    pqc::VulnDatabase db;
    db.load("data/vulnerable_functions.json");

    pqc::ProjectScanner scanner;
    auto inv = scanner.scan("tests/fixtures"); // инвентаризация

    pqc::RegexAnalyzer ra(db);
    auto findings = ra.analyze(inv);           // анализ всех SOURCE-файлов
    ASSERT_NOT_EMPTY(findings);

    pqc::RiskAssessor assessor(db);
    auto rr = assessor.assess(findings, inv);  // оценка риска и план миграции

    ASSERT_NOT_EMPTY(rr.migration_plan);       // план не пуст
    ASSERT_GT(rr.overall_risk_score, 0.0);     // итоговый риск рассчитан
    ASSERT_NOT_EMPTY(rr.nist_readiness);       // статус готовности к NIST определён
}
\end{itmolisting}

Тест \texttt{regex\_precision\_recall\_gt50pct}
(листинг~\ref{lst:itest-metrics}) является регрессионным:
он загружает эталонную разметку и проверяет, что все три
метрики качества не опускаются ниже~50\,\%.
Порог выбран консервативным, чтобы тест оставался стабильным
при расширении базы, но немедленно сигнализировал о грубом
регрессе.

\begin{itmolisting}[C++]{Интеграционный регрессионный тест метрик качества}{lst:itest-metrics}
TEST(FullPipeline, regex_precision_recall_gt50pct) {
    pqc::VulnDatabase db;
    db.load("data/vulnerable_functions.json");
    pqc::ProjectScanner scanner;
    auto inv = scanner.scan("tests/fixtures");
    pqc::RegexAnalyzer ra(db);
    auto findings = ra.analyze(inv);

    // Загружаем эталонную разметку из ground_truth.json (32 записи)
    auto gt = pqc::MetricsCalculator::load_from_json(
        "tests/fixtures/ground_truth.json"
    );
    pqc::MetricsCalculator calc;
    auto m = calc.compute(findings, gt);

    // Минимальные пороги: даже при частичном покрытии базы
    // инструмент должен демонстрировать базовое качество обнаружения.
    ASSERT_GE(m.precision(), 0.5);
    ASSERT_GE(m.recall(),    0.5);
    ASSERT_GE(m.f1_score(),  0.5);
}
\end{itmolisting}

Тест \texttt{migration\_plan\_has\_fips\_references} проверяет
сквозное прохождение ссылок на стандарты: поле
\texttt{replacement\_standard} хотя бы одной записи плана
миграции должно содержать \texttt{"FIPS"}~--- это подтверждает
корректную загрузку и применение рекомендаций из базы.
Тест \texttt{risk\_report\_has\_tc26\_data} аналогично проверяет
наличие непустого поля \texttt{tc26\_note}, что критично для
применимости инструмента в проектах с ГОСТ"=криптографией.

% -----------------------------------------------------------
\subsection{Оценка результатов тестирования}

Все~47 тестов проходят успешно.
Итоги приведены в таблице~\ref{tab:test_results_summary}.

\begin{longtable}{|p{5.5cm}|r|r|}
  \caption{Итоги выполнения тестов}
  \label{tab:test_results_summary} \\
  \hline
  Набор тестов & Всего & Пройдено \\ \hline
  \endfirsthead
  \multicolumn{3}{r}{\textit{Продолжение таблицы \thetable}} \\ \hline
  Набор тестов & Всего & Пройдено \\ \hline
  \endhead \hline \endlastfoot

  \texttt{VulnDatabase}               & 9  & 9  \\ \hline
  \texttt{ProjectScanner}             & 5  & 5  \\ \hline
  \texttt{RegexAnalyzer}              & 8  & 8  \\ \hline
  \texttt{ASTAnalyzer}                & 7  & 7  \\ \hline
  \texttt{RiskAssessor}               & 7  & 7  \\ \hline
  \texttt{MetricsCalculator}          & 5  & 5  \\ \hline
  \texttt{FullPipeline} (интеграционные) & 6 & 6 \\ \hline
  \textbf{Итого}                      & \textbf{47} & \textbf{47} \\

\end{longtable}

Интеграционный тест \texttt{regex\_precision\_recall\_gt50pct}
зафиксировал финальные метрики качества на эталонном наборе из
32~записей: для regex"=режима~--- precision\,=\,80,0\,\%,
recall\,=\,100,0\,\%, F1\,=\,88,9\,\%; для AST"=режима~---
precision\,=\,94,1\,\%, recall\,=\,100,0\,\%, F1\,=\,97,0\,\%.
Детальный анализ приведён в разделе~\ref{sec:analysis_comparison}.


% ============================================================
%  АНАЛИЗ ПОЛУЧЕННЫХ РЕЗУЛЬТАТОВ
% ============================================================

\section{Анализ полученных результатов}

\subsection{Соответствие реализации поставленным требованиям}

В таблице~\ref{tab:req_compliance} приведено соответствие
реализованных компонентов ключевым функциональным требованиям.

\begin{longtable}{|p{5.8cm}|p{4.6cm}|p{4.6cm}|}
  \caption{Соответствие реализации функциональным требованиям}
  \label{tab:req_compliance} \\
  \hline
  Требование & Реализация & Подтверждение \\ \hline
  \endfirsthead
  \multicolumn{3}{r}{\textit{Продолжение таблицы \thetable}} \\ \hline
  Требование & Реализация & Подтверждение \\ \hline
  \endhead \hline \endlastfoot

  Обнаружение квантово"=уязвимых функций в C/C++"=коде
    & Два анализатора: \texttt{RegexAnalyzer} и \texttt{ASTAnalyzer}
    & TP\,=\,32/32 в обоих режимах \\ \hline

  Поддержка OpenSSL, Crypto++, ГОСТ"=крипто\-графии
    & Три секции в \texttt{vulnerable\_functions.json}
    & Тест \texttt{gost\_library\_entries} \\ \hline

  Условная уязвимость на основе аргументов вызова
    & \texttt{ASTAnalyzer}: \texttt{resolve\_expr} + PQC"=фильтр
    & FP\,=\,0 на \texttt{safe\_pqc\_kem} в AST"=режиме \\ \hline

  Расчёт риска с учётом контекста
    & \texttt{RiskAssessor}: 6~контекстных коэффициентов
    & Тесты \texttt{network\_file\_increases\_risk},
      \texttt{test\_file\_reduces\_risk} \\ \hline

  Приоритизированный план миграции
    & \texttt{RiskAssessor::assess}: сортировка по \texttt{final\_score}
    & Тест \texttt{migration\_order\_sorted\_by\_risk} \\ \hline

  Рекомендации с ссылками FIPS и ТК\,26
    & Поля \texttt{replacement\_standard}, \texttt{tc26\_note}
      в \texttt{MigrationAction}
    & Тесты \texttt{migration\_plan\_has\_fips\_references},
      \texttt{risk\_report\_has\_tc26\_data} \\ \hline

  CBOM в формате CycloneDX
    & Класс \texttt{Cbom}, метод
      \texttt{ReportGenerator::\allowbreak build\_cbom}
    & Тест \texttt{report\_json\_serialization} \\ \hline

  Расширяемость пользовательской базой
    & \texttt{VulnDatabase::merge\_json}
    & Тест \texttt{merge\_extra\_db} \\ \hline

  Измерение precision/recall/F1
    & \texttt{MetricsCalculator}, флаг \texttt{{-}{-}metrics}
    & Тест \texttt{regex\_precision\_recall\_gt50pct} \\

\end{longtable}

% -----------------------------------------------------------
\subsection{Сравнение режимов анализа}
\label{sec:analysis_comparison}

Оба режима тестировались на файле \texttt{sample\_vulnerable.cpp}
(32 записи в эталонной разметке).
Результаты представлены в таблице~\ref{tab:mode_comparison}.

\begin{longtable}{|p{4.8cm}|r|r|}
  \caption{Сравнение метрик качества AST и regex"=режимов}
  \label{tab:mode_comparison} \\
  \hline
  Метрика & AST & Regex \\ \hline
  \endfirsthead
  \multicolumn{3}{r}{\textit{Продолжение таблицы \thetable}} \\ \hline
  Метрика & AST & Regex \\ \hline
  \endhead \hline \endlastfoot

  Эталонных записей (GT)     & 32       & 32       \\ \hline
  True Positives (TP)         & 32       & 32       \\ \hline
  False Positives (FP)        & 2        & 8        \\ \hline
  False Negatives (FN)        & 0        & 0        \\ \hline
  Precision                   & 94,1\,\% & 80,0\,\% \\ \hline
  Recall                      & 100,0\,\%& 100,0\,\%\\ \hline
  F1                          & 97,0\,\% & 88,9\,\% \\ \hline
  Время анализа тестового файла
    & $\approx$1030\,мс & $\approx$80\,мс \\

\end{longtable}

\paragraph{True Positives.}
Оба режима обнаружили все~32 ожидаемые находки.
Это охватывает широкий спектр квантово"=уязвимых конструкций:
прямые вызовы низкоуровневых функций (\texttt{RSA\_new},
\texttt{ECDSA\_sign}, \texttt{DH\_new}),
EVP"=функции высокого уровня с классическим аргументом
(\texttt{EVP\_PKEY\_CTX\_new\_from\_name("RSA",...)},
полные цепочки подписи и верификации через
\texttt{EVP\_DigestSign*\,/\,EVP\_DigestVerify*},
хеш"=функции (\texttt{SHA1}, \texttt{MD5}),
а также функции, доступные только через псевдонимы в базе
(\texttt{EC\_KEY\_generate\_key},
\texttt{DH\_generate\_parameters\_ex}).

\paragraph{False Positives.}
\textbf{AST (2~FP).}
Обе находки~--- \texttt{BN\_new} на строках~50 и~217~---
являются фактически корректными срабатываниями: создание
BIGNUM"=структуры служит надёжным индикатором ручной реализации
RSA или DH (квантово"=уязвимых алгоритмов), но эти строки не
вошли в эталонную разметку.
Если включить их в GT, precision AST достигает~100\,\%.

\textbf{Regex (8~FP).}
Помимо тех же двух \texttt{BN\_new}~--- 6~находок из блока
\texttt{SAFE}: три вызова в функции \texttt{safe\_pqc\_kem}
(\texttt{EVP\_PKEY\_CTX\_new\_from\_name("ML-KEM-768",...)},
\texttt{EVP\_PKEY\_keygen\_init}, \texttt{EVP\_PKEY\_keygen})
и аналогичные три вызова из второго PQC"=безопасного контекста.
Regex"=анализатор не может отследить провенанс аргумента
\texttt{"ML-KEM-768"}, поэтому срабатывает на имя функции
вне зависимости от переданного алгоритма.
AST"=анализатор корректно подавляет все эти~6~находок через
механизм \texttt{resolve\_expr}: при обнаружении постквантового
имени алгоритма в цепочке аргументов находка не формируется.

\paragraph{False Negatives.}
В обоих режимах ложноотрицательных результатов нет.
Полнота достигается за счёт покрытия базы: она включает
как основные имена функций, так и их псевдонимы, а также
отдельные записи для хеш"=функций и функций генерации параметров,
которые иначе могли бы быть пропущены.

\paragraph{Рекомендации по выбору режима.}
Regex"=режим целесообразно использовать как быстрый первичный
скрининг: время обработки на тестовом файле~--- около~80\,мс
против~1030\,мс у AST"=режима.
AST"=режим рекомендуется как основной для формирования итогового
отчёта, особенно в проектах, активно использующих EVP"=API
с передачей алгоритма через строковый аргумент.

% -----------------------------------------------------------
\subsection{Оценка практической применимости}

Помимо поиска уязвимостей, инструмент формирует три класса
выходных данных, пригодных для аудита и регуляторной отчётности:
\begin{itemize}
  \item \textbf{CBOM"=файл} в формате CycloneDX~--- стандартизованный
    инвентарь криптографических активов с полями
    \texttt{primitive} (rsa, ecdh, sha1~и~др.),
    \texttt{cryptoFunctions} (keyGeneration, keyEncapsulation,~\dots)
    и \texttt{quantumSecurityLevel} (0~для полностью уязвимых,
    50~для частично ослабленных, 128~для квантово"=устойчивых);
  \item \textbf{план миграции} с явными ссылками на FIPS\,203/204/205
    и рекомендации ТК\,26 для переходного периода 2025--2030, включая
    примеры кода для каждой замены;
  \item \textbf{оценку NIST"=готовности} (\texttt{not-started\,/\,
    in-progress\,/\,near-compliant\,/\,compliant}), вычисляемую
    на основе числа критических и высоких находок.
\end{itemize}
Время AST"=анализа файла объёмом~$\approx$2000~строк составляет
менее~2~с~--- это приемлемо для встраивания в CI"=конвейер.

% -----------------------------------------------------------
\subsection{Ограничения реализации}

\paragraph{Внутрипроцедурный анализ аргументов.}
Механизм \texttt{resolve\_expr} прослеживает провенанс аргументов
только в рамках одной функции.
Если значение алгоритма передаётся через параметр функции или
хранится в поле класса, инициализированном в другом методе,
анализатор не сможет определить постквантовую природу аргумента
и сформирует ложное срабатывание.

\paragraph{Провенанс аргументов в режиме regex.}
Regex"=анализатор не отслеживает значения переменных, что
структурно приводит к~6 неустранимым FP для EVP"=функций
с постквантовыми аргументами.
Любые эвристические правила (белые списки имён переменных)
будут хрупкими и не масштабируются на произвольный код.

\paragraph{Зависимость от libclang.}
AST"=анализатор требует установленной libclang и корректно
настроенных путей включения.
При неполной конфигурации заголовочных файлов часть вызовов
может не разрешиться, и анализатор переключается на
токенизирующий fallback"=режим с пониженной точностью
извлечения контекста.

\paragraph{Покрытие базы данных.}
Текущая база охватывает три библиотеки: OpenSSL, Crypto++
и ГОСТ"=криптографию через OpenSSL"=провайдер.
Функции из wolfSSL, Botan, mbedTLS и других криптографических
библиотек не включены и не будут обнаружены.

\paragraph{Отсутствие модуля автоматической миграции.}
Инструмент предоставляет готовый план миграции с примерами кода,
однако не применяет кодовые преобразования автоматически~---
разработчик выполняет замену вручную на основе сгенерированного
отчёта.

% -----------------------------------------------------------
\subsection{Направления дальнейшего развития}

\paragraph{Модуль автоматической миграции кода.}
Приоритетным направлением является реализация
автоматического трансформатора AST, который по записям плана
миграции и шаблонам из базы данных генерирует готовые патчи.
Например, для нахождения \texttt{RSA\_generate\_key\_ex}
трансформатор должен сформировать эквивалентный блок с
\texttt{EVP\_PKEY\_CTX\_new\_from\_name("ML-DSA-65",~...)},
сохраняя семантику вызова.
Технически это требует расширения AST"=обходчика до уровня
переписывающего трансформера на базе того же cursor"=механизма
libclang.

\paragraph{Межпроцедурный анализ аргументов.}
Реализация хотя бы однопроходного межпроцедурного анализа
в рамках одной единицы трансляции позволит корректно обрабатывать
передачу алгоритма через параметры функций~--- характерный
паттерн в проектах с криптографическими фасадами.
Это снизит число как ложноположительных, так и
ложноотрицательных результатов на реальном коде.

\paragraph{Расширение базы данных.}
Добавление функций из wolfSSL, mbedTLS, Botan, а также
Python"=биндингов (\texttt{cryptography}, \texttt{pyOpenSSL}),
Go"=стандартной библиотеки (\texttt{crypto/rsa},
\texttt{crypto/elliptic}) и Rust"=крейтов (\texttt{rsa},
\texttt{p256}) сделает инструмент применимым для
многоязычных проектов.

\paragraph{Интеграция в CI/CD.}
Добавление режима выхода с ненулевым кодом при обнаружении
критических находок (\texttt{{-}{-}fail"=on"=critical}) и
формирование SARIF"=отчёта позволят встраивать инструмент
в GitHub~Actions, GitLab~CI и аналогичные системы в качестве
обязательного шага перед слиянием.

\paragraph{Временно"=зависимая оценка риска.}
Текущий оценщик риска не учитывает временной горизонт
появления криптографически значимых квантовых компьютеров.
Добавление временно"=зависимой компоненты~--- например, повышающего
коэффициента для долгосрочных ключей с учётом прогноза к~2030~г.~---
сделает приоритизацию плана миграции более реалистичной и
практически значимой.
